%======================================================================
\chapter{Theorem 10}\label{ch:thm10}
%======================================================================

\begin{theorem}[Bing and Jones, Theorem~10]\label{thm:bj10}
Let $X$, $Y$ be snake-like compact metric spaces and $f\colon X\to\I$, $g\colon Y\to\I$ continuous arcs of pseudo-arcs. Let $h_0\colon\fib f0\to\fib g0$ and
$h_1\colon\fib f1\to\fib g1$ be homeomorphisms. There is a homeomorphism $H\colon X\to Y$ with $g\circ H=f$ which extends $h_0$ and $h_1$.
\end{theorem}

\formal{\lean{HPC.BJ59.bjThm10}, a proof of the statement \lean{HPC.BJ59.BJThm10}.}

The homeomorphism is produced by Theorem~\ref{thm:zigzag}, from chains built by one construction (the \emph{step}) made alternately on the two sides. We write
$p_x=\fib fx$, $q_x=\fib gx$, $P(J)=\fib fJ$ and $Q(J)=\fib gJ$. A family of links \emph{properly covers} a set if each of its links meets the set and their union
contains it.

\section{The conditions carried along}

\begin{definition}[Bing and Jones, p.~188]\label{def:conditions}
Let $D$ be a chain on $X$, $E$ a chain on $Y$, and $H$ a map of the indices of $D$ to those of $E$.
\begin{description}
\item[\textnormal{(A)}] $H$ is a chain map of $D$ onto $E$; if $D_j$ meets $p_x$ then $E_{H(j)}$ meets $q_x$; and if a run of $D$ properly covers $p_x$, the images
of its links cover $q_x$.
\item[\textnormal{(C)}] If $x\in p_0\cap D_j$ then $h_0(x)\in E_{H(j)}$; if $x\in p_1\cap D_j$ then $h_1(x)\in E_{H(j)}$.
\end{description}
\end{definition}

\formal{\lean{HPC.BJ59.IsFibreChainMap} (Condition A), \lean{HPC.BJ59.EndCompat} (Condition C).}

Recall private points (Definition~\ref{def:private}) and fibre-fineness (Definition~\ref{def:fibre-fine}). We also need a name for the conclusion of the
corollary to Theorem~8: a map $R$ of the indices of a chain $F$ \emph{runs through $[\sigma,\sigma']$ on the run $[t,u]$} if every $j\in[\sigma,\sigma']$ is the value of $R$ at
three consecutive indices of $[t,u]$.

\formal{\lean{HPC.BJ59.IsPatternOn}.}

\section{The step}

\begin{theorem}[The step]\label{thm:step}
Let $X$, $Y$, $f$, $g$, $h_0$, $h_1$ be as in Theorem~\ref{thm:bj10}. Let $D$ be a chain covering $X$ with private points, $E$ a chain covering $Y$ which is
fibre-fine for $g$, and $H$ satisfying (A) and (C). For every $\delta>0$ there are a BJ-chain $E'$ covering $Y$, of mesh $<\delta$, with private points, and a map
$K$ of its indices to those of $D$ such that:
\begin{enumerate}[label=(\roman*)]
\item (A$'$) $K$ satisfies (A) with the roles of $X$ and $Y$ exchanged;
\item (C$'$) $K$ satisfies (C) for $h_0^{-1}$, $h_1^{-1}$;
\item (B$'$) $E'_k\subseteq E_{H(K(k))}$ for every $k$.
\end{enumerate}
\end{theorem}

\formal{The statement is \lean{HPC.BJ59.BJStep}, proved as \lean{HPC.BJ59.bj_step}.}

Applied with $X$ and $Y$ exchanged (to $E'$, $D$, $K$, $h_0^{-1}$, $h_1^{-1}$), the step returns a chain $D'$ on $X$ and a chain map $H'$ of $D'$ onto $E'$
with (A), (C) and $D'_j\subseteq D_{K(H'(j))}$, which is Condition~(B) of Chapter~\ref{ch:zigzag}. So one statement serves both halves of the induction.

The construction follows Bing and Jones's Steps~2 and~3 (pp.~183--188), with two differences: the homeomorphisms $h_0$, $h_1$ are used only on the end fibres,
and at an interior cut the construction passes from the map of one sub-arc to the map of the next through an overlap of the sub-arcs. We first assemble the
pieces.

\begin{lemma}[The map at an end; Step~2(b)]\label{lem:end-piece}
Let $c\in\set{0,1}$ and let $h\colon p_c\to q_c$ be a homeomorphism. Let $W=[w_1,w_2]$ be a run of $D$ covering $p_c$ from which neither end link can be dropped. There
is $\varepsilon>0$ such that for every $\varepsilon$-chain $F$ covering $Y$ there is a chain map $R$ of $F$ into $W$ with
\begin{itemize}
\item $h^{-1}(\cl{F_a}\cap q_c)\subseteq D_{R(a)}$ for every link $F_a$;
\item every $j\in W$ is the value of $R$ at three consecutive links meeting $q_c$.
\end{itemize}
\end{lemma}

\begin{proof}
By Lemma~\ref{lem:runs}(c) every $j\in W$ has a point $\varphi(z)$, $z\in q_c$, $\varphi=h^{-1}$, lying in $D_j$ and in no other link of $W$. Let $\lambda$ be a Lebesgue number of the
cover of $p_c$ by the links of $W$, and choose $\varepsilon$ below a third of $\diam q_c$ and such that $\varphi$ maps subsets of $q_c$ of diameter $<5\varepsilon$ to sets of
diameter $<\lambda$. Let $T$ be the run of links of $F$ meeting $q_c$. For $a\in T$ let $S_a$ be the set of $j\in W$ with $\varphi(\cl{F_a}\cap q_c)\subseteq D_j$: a run of one
or two indices, nonempty. Call $a$ \emph{decided} if $S_a=\set j$. If $a$ is decided with value $j$, then $j\in S_{a'}$ for every $a'\in T$ within distance $4$ of $a$ (the
closures of the links from $a$ to $a'$ meet $q_c$ in a set of diameter $<5\varepsilon$, whose image lies in one link of $W$, necessarily $D_j$); so decided positions at most four
apart have the same value. Put $R(a)=j$ if some decided position within distance $2$ of $a$ has value $j$, and $R(a)=\min S_a$ otherwise; extend $R$ outside $T$ by its values at the
ends of $T$. Then $R(a)\in S_a$, and $R$ is a chain map. For $j\in W$, take $z$ as above and $a\in T$ with $z\in F_a$: $a$ is decided with value $j$, so $R=j$ at the
positions of $T$ within distance $2$ of $a$, at least three of them.
\end{proof}

\formal{\lean{HPC.BJ59.exists_end_piece}, from \lean{HPC.BJ59.exists_pattern}.}

\begin{lemma}[The map of a sub-arc]\label{lem:subarc-piece}
Let $a<b$ and suppose every run of $D$ covering $P[a,b]$ minimally has all its links meeting every fibre over $[a,b]$ (Theorem~\ref{thm:step1b}). Then there is a run
$[r,s]$ of $D$ properly covering every $p_x$, $x\in[a,b]$, and $\varepsilon>0$ such that for every $\varepsilon$-chain $F$ covering $Y$ there is a chain map $R$ of $F$ into $[r,s]$
with
\begin{description}
\item[\textnormal{(R1)}] $F_k\cap Q[a,b]\subseteq E_{H(R(k))}$ for every $k$;
\item[\textnormal{(R2)}] for $x\in[a,b]$, $R$ runs through $[r,s]$ on every run of $F$ covering $q_x$.
\end{description}
\end{lemma}

\begin{proof}
Take a minimal covering run $[r,s]$ of $P[a,b]$. By (A), its image $H([r,s])$ (a run, Lemma~\ref{lem:chainmap-ivt}) consists of links meeting every $q_x$, $x\in[a,b]$, and covers
$Q[a,b]$. By fibre-fineness it has one link or at least three. Apply the corollary to Theorem~8 on the sub-arc $Q[a,b]$ (Corollary~\ref{cor:bj8} and the remark after
it) with this run as $E$ and $H$ restricted to $[r,s]$ as $H_1$.
\end{proof}

\formal{\lean{HPC.BJ59.exists_subarc_piece}.}

\begin{lemma}[Stability near a fibre]\label{lem:tau}
If a run of a chain properly covers $\fib\varphi c$ for a continuous open $\varphi$ onto $\I$, it properly covers $\fib\varphi x$ for all $x$ near $c$. If the closure of
every link $F_k$ meets $\fib\varphi c$ inside an open set $V_k$, there is $\tau>0$ with $F_k\cap\fib\varphi{[c-\tau,c+\tau]}\subseteq V_k$ for all $k$.
\end{lemma}

\formal{\lean{HPC.BJ59.exists_tau_run}, \lean{HPC.BJ59.exists_tau_incl}.}

\begin{proof}[Proof of Theorem~\ref{thm:step} in outline]
\emph{(a) Sub-arcs.} Let $\delta_D$ be given by Theorem~\ref{thm:step1b} for $f$ and $D$. Choose $N_0\ge1$ and $\eta>0$ with $1/N_0+2\eta\le\delta_D$ and $\eta<1/(4N_0)$. Put
$c_i=i/N_0$ and $J_i=[c_i-\eta,c_{i+1}+\eta]\cap\I$ for $i<N_0$; consecutive $J_i$ overlap around $c_{i+1}$.

\emph{(b), (c) The pieces.} Lemma~\ref{lem:subarc-piece} on each $J_i$ gives runs $[r_i,s_i]$ and numbers $\varepsilon_i$; Lemma~\ref{lem:end-piece} at $c=0$, $1$ (with $h=h_0$,
$h_1$) gives runs $W^0$, $W^1$ and numbers $\varepsilon^0$, $\varepsilon^1$.

\emph{(d) The chain $F$.} Let $F$ be an $\varepsilon$-chain covering $Y$, with $\varepsilon$ below all $\varepsilon_i$, $\varepsilon^c$ and so small that $g$ maps each link of $F$
to a set of diameter $<\eta/4$. We obtain chain maps $R_i\colon F\to[r_i,s_i]$ with (R1), (R2) on $J_i$, and $R^c\colon F\to W^c$ as in Lemma~\ref{lem:end-piece}.

\emph{(e) Windows.} Let $T^0$, $T^1$ be the runs of $F$ meeting $q_0$, $q_1$, and $T_i$ ($0<i<N_0$) the run meeting $q_{c_i}$. By Lemma~\ref{lem:tau} choose $\tau\in(0,\eta/4)$ such
that these runs properly cover the fibres within $\tau$ of $0$, $1$, $c_i$; $W^c$ properly covers $p_x$ for $\abs{x-c}\le\tau$; and $F_a\cap Q(\set{\abs{x-c}\le\tau})\subseteq
E_{H(R^c(a))}$ for $F_a\in T^c$ (at $x=c$ this follows from Lemma~\ref{lem:end-piece} and (C)). The windows are $L^0=[\tau/4,\tau/2]$, $L^1=[1-\tau/2,1-\tau/4]$ and
$L_i=[c_i,c_i+\tau/2]$. At each window two maps meet: $R^0$ and $R_0$ at $L^0$, $R_{i-1}$ and $R_i$ at $L_i$, $R_{N_0-1}$ and $R^1$ at $L^1$. Both carry the run $T$ of
that window onto runs of $D$ properly covering the fibre at the centre, so Theorem~\ref{thm:bj6} gives a \emph{transition link} $F_y\in T$ at which their values are adjacent.
Lemma~\ref{lem:window-set} gives a transition set $U\subseteq F_y$ over the interior of the window, and Theorem~\ref{thm:bj5-window} splits $Y\setminus U$ into closed sides
$A$ (left) and $B$ (right) such that, for each $x$ in the window, every link of $T$ meets $A\cap q_x$ or every link meets $B\cap q_x$. Intersecting consecutive sides gives
closed, pairwise disjoint regions
\[
Z^0=A^0,\qquad Z_i=B_i\cap A_{i+1}\ (i<N_0,\ B_0:=B^0,\ A_{N_0}:=A^1),\qquad Z^1=B^1,
\]
covering the complement of the transition sets, with $Z_i\subseteq Q(J_i)$; region $Z_i$ carries the map $R_i$, and $Z^c$ the map $R^c$.

\emph{(f) $E'$ and $K$.} Choose $\mu>0$ below $\delta$, a Lebesgue number of $F$, the distances between distinct regions, between each transition set and the regions not next
to it, between the two sides of each window and from each transition set to the complement of its link $F_y$, and so small that a set of diameter $<\mu$ has image under $g$
of diameter $<\tau/4$. Then a set of diameter $<\mu$ meets at most one transition set, and none if it meets $q_0\cup q_1$. Let $G$ be a BJ-chain of mesh $<\mu$ covering $Y$
(Theorem~\ref{thm:shrink}). For a link $e$ of $G$: if $e$ meets a transition set $U\subseteq F_y$, then $e\subseteq F_y$, and $K(e)$ is the left map at $y$ if $e$ meets the side $A$,
otherwise the right map at $y$; if not, $e$ meets exactly one region, and $K(e)$ is that region's map at a link of $F$ containing $e$.

\emph{(g) Verification.}
\begin{enumerate}
\item $K$ is a chain map: inside a region adjacent links of $G$ lie in equal or adjacent links of $F$; a link meeting $U$ and an adjacent link not meeting it meet on the same side
and use the same map; two links meeting $U$ lie in $F_y$, where the two maps are adjacent.
\item (B$'$) follows from (R1) in the regions $Z_i$ and at interior transitions (where $e\subseteq F_y\cap Q(J_{i-1}\cap J_i)$), and from the choice of $\tau$ in $Z^c$.
\item (C$'$): a link meeting $q_c$ meets no transition set, so it lies in $Z^c$, and Lemma~\ref{lem:end-piece} applies.
\item (A$'$), first clause: if $e$ meets $q_x$, $x$ lies in the range of the map used, whose values all name links meeting $p_x$.
\item (A$'$), second clause. Let a run $C$ of $G$ properly cover $q_x$. We find a run $S$ of $D$ properly covering $p_x$, a map $M$ among the $R_i$, $R^c$, a run $T$ of $F$
and a set $\Gamma\subseteq q_x$ meeting every link of $T$, such that every $j\in S$ is the value of $M$ at three consecutive links $F_{a-1},F_a,F_{a+1}$ of $T$ and every
link of $G$ meeting $\Gamma$ takes its value from $M$. Then for $j\in S$, a point $z\in\Gamma\cap F_a$ lies in some $e\in C$, the link of $F$ used for $e$ contains $z$, so it is
$F_{a-1}$, $F_a$ or $F_{a+1}$, and $K(e)=j$. Between windows, $q_x$ lies in one region $Z_i$: take $M=R_i$, $\Gamma=q_x$, $S=[r_i,s_i]$ and (R2). In a window, the dichotomy of
Theorem~\ref{thm:bj5-window} says that every link of $T$ meets $A\cap q_x$ (take the left map, $\Gamma=A\cap q_x$) or every link meets $B\cap q_x$ (the right map); (R2) at the
centre of the window, or Lemma~\ref{lem:end-piece} at the ends, gives the three consecutive links. Near $0$ and $1$ use $R^c$ and $W^c$.
\item $K$ is onto: for a private point $z$ of $D_j$ and $x=f(z)$, the run $S$ of item~5 covers $z$, so $j\in S$ is a value of $K$.
\item Private points: shrink $G$ by Proposition~\ref{prop:private-shrink} and keep $K$. Every property survives, since links only shrink and a run of shrunk links properly
covering $q_x$ comes from a run of $G$ properly covering it.\qedhere
\end{enumerate}
\end{proof}

\formal{The gluing (e)--(g) is \lean{HPC.BJ59.bj_glue}; the pieces are \lean{HPC.BJ59.exists_subarc_piece}, \lean{HPC.BJ59.exists_end_piece}; the windows use
\lean{HPC.BJ59.adjacent_of_properCover}, \lean{HPC.BJ59.exists_window_set}, \lean{HPC.BJ59.bj_theorem5_window}.}

\section{The induction}

\begin{proof}[Proof of Theorem~\ref{thm:bj10}]
By Lemma~\ref{lem:fibre-spread} let $\kappa_X$, $\kappa_Y$ bound the spreads of the fibres of $f$ and $g$ from below. A \emph{stage} is a triple $(D,E,H)$ of chains covering
$X$ and $Y$, BJ-chains with private points, $D$ fibre-fine for $f$ and $E$ for $g$, and $H$ satisfying (A) and (C), with both meshes below a given $\varepsilon$.

\emph{Stage $0$}: the one-link chains $(X)$, $(Y)$ and the identity map. All conditions hold trivially (a one-link chain is fibre-fine, its link has private points, (A) and
(C) hold because the only link of $E$ is $Y$).

\emph{From stage $i$ to stage $i+1$}: apply Theorem~\ref{thm:step} to $(D^i,E^i,H_i)$ with $\delta_i=\min(2^{-i-1},\kappa_X/3,\kappa_Y/3)$, obtaining $E^{i+1}$ and $K_i$;
then apply it with $X$ and $Y$ exchanged to $(E^{i+1},D^i,K_i)$, obtaining $D^{i+1}$ and $H_{i+1}$. By Lemma~\ref{lem:fibre-fine} the new chains are fibre-fine. Condition
(B$'$) of Chapter~\ref{ch:zigzag} is (iii) of the first application and (B) is (iii) of the second.

\emph{The homeomorphism.} Theorem~\ref{thm:zigzag} with $P=X$, $Q=Y$ gives a homeomorphism $H\colon X\to Y$ with $H(p)\in\cl{\st_i(p)}$ for every $i$.
Proposition~\ref{prop:zigzag-fibre}, whose hypothesis is the first clause of (A), gives $g\circ H=f$. For $x\in p_0$, Condition (C) gives $h_0(x)\in\st_i(x)$ for every $i$, so
$h_0(x)=H(x)$ by the characterization; likewise for $h_1$.
\end{proof}

\formal{\lean{HPC.BJ59.bjThm10_of_step}, with \lean{HPC.BJ59.BJStage}, \lean{HPC.BJ59.BJStage.exists_init}, \lean{HPC.BJ59.BJStage.exists_next}.}

\begin{remark}
Bing and Jones extend $h$ at every step to the fibres over the cut points $i/n$ (arbitrarily in Step~1(c), by their Theorem~9 in Step~2(i)) and build the maps near those
fibres from it. In the proof above $h_0$ and $h_1$ are used only on $p_0$ and $p_1$, and Theorem~9 is not needed. There is also no separate first step: the induction starts
from one-link chains, and for a one-link chain the corollary to Theorem~8 is the zigzag map of their Step~2(b).
\end{remark}

\formal{Bing and Jones's Theorem~9 is nevertheless proved: \lean{HPC.BJ59.bj_theorem9}.}

\section*{Exercises}

\begin{exercise}
Check that the one-link stage satisfies (A): what does ``a run of $D$ properly covering $p_x$'' mean when $D$ has one link?
\end{exercise}

\begin{exercise}
In Lemma~\ref{lem:end-piece}, show that two runs $S_a$, $S_{a+1}$ of length at most two that meet have least elements at most one apart, and deduce that $R$ is a chain map.
\end{exercise}

\begin{exercise}
Show that Theorem~\ref{thm:bj10} gives, for any homeomorphism $h_0\colon p_0\to p_0$ of the end fibre of a continuous snake-like arc of pseudo-arcs, a self-homeomorphism of $X$
preserving every fibre and extending $h_0$.
\end{exercise}
